\documentclass[10pt,a4paper]{amsart}
\usepackage[margin=19mm]{geometry}
\usepackage{amsmath,amssymb,mathtools}
\usepackage[scr=rsfs,cal=euler]{mathalfa}
\usepackage{xfrac}
\usepackage[unicode,hidelinks,bookmarks=false]{hyperref}
\newcommand{\ie}{i.e.\ }
\newcommand{\cf}{cf.\ }
\newcommand{\resp}{respectively\ }
\newcommand{\Subsection}{\subsection}
\providecommand{\nobreakdash}{\nobreak\hbox{-}}
\setlength{\emergencystretch}{2em}
\multlinegap0pt
\usepackage{mathtools}
\usepackage{enumitem}
\usepackage{cleveref}
\newtheorem{theorem}{Theorem}[section]
\newtheorem{lemma}[theorem]{Lemma}
\newtheorem{proposition}[theorem]{Proposition}
\newtheorem{corollary}[theorem]{Corollary}
\newtheorem{claim}[theorem]{Claim}
\newtheorem{definition}[theorem]{Definition}
\newtheorem{remark}[theorem]{Remark}
\newcommand{\N}{\mathbb{N}}
\newcommand{\Z}{\mathbb{Z}}
\newcommand{\one}{\mathbf{1}}
\newcommand{\Oh}{O(1)}
\newcommand{\interval}[2]{[ #1,#2 ]}
\begin{document}
\title[A sharp 5/8 bound for an ErdHos-S'os pairwise-sums problem]{A sharp 5/8 bound for an Erd\H{o}s-S\'os pairwise-sums problem}
\author{Ricky Cipollini}
\date{}
\maketitle
\noindent\emph{Source and licence.} A sharp 5/8 bound for an Erd\H{o}s-S\'os pairwise-sums problem. Version arXiv:2606.29361v1 (CC BY 4.0). This material is distributed under the Creative Commons Attribution 4.0 International licence, http://creativecommons.org/licenses/by/4.0/, as recorded in the arXiv OAI licence metadata for this exact version and shown on the abstract page https://arxiv.org/abs/2606.29361v1. Full rights notice: licenses/arxiv-2606.29361-CC-BY-4.0.txt.\par\smallskip
\noindent\textit{Editorial scope.} Counted unit: the original \texttt{lemma} environment \texttt{lem:folding} at source lines 259--268 together with its complete proof at lines 270--293; the delivered document keeps source lines 90--294 so that every referenced label and every citation resolves inside it. Native editable source is the paper's own concatenated TeX; the AMS wrapper is editorial formatting only. No publisher class, upstream script or unrelated artwork is redistributed.
\begin{equation}\label{eq:foldedOK}
  x+y\neq m,
  \qquad
  x+y\notin B\pmod m.
\end{equation}
Define
\[
  S_0(B)=\{x+y:x\neq y,\ x,y\in B,\ x+y<m\},
\]
\[
  S_1(B)=\{x+y-m:x\neq y,\ x,y\in B,\ x+y>m\},
\]
and
\[
  C(B)=S_0(B)\cap S_1(B).
\]
Thus $C(B)$ is the set of residues which occur both as a non-wrapped and as a wrapped pair-sum of two distinct elements of $B$.

\begin{lemma}[Folded additive lemma]\label{lem:folded}
For every $m$ and $B$ satisfying \eqref{eq:foldedOK},
\[
  |B|-|C(B)|\leq \frac m4+2.
\]
\end{lemma}

\begin{proof}
We induct on $|B|$.

First observe the reflection symmetry. Put
\[
  -B=\{m-b:b\in B\}\subseteq \{1,\ldots,m-1\}.
\]
Replacing $B$ by $-B$ preserves \eqref{eq:foldedOK}; low sums for $B$ become high sums for $-B$, and high sums for $B$ become low sums for $-B$. Hence
\begin{equation}\label{eq:collision-reflection}
  |C(-B)|=|C(B)|.
\end{equation}

Assume $|B|\geq 2$. After replacing $B$ by $-B$ if necessary, we may suppose that the two smallest elements of $B$,
\[
  \alpha<\beta,
\]
satisfy
\begin{equation}\label{eq:small-sum}
  \alpha+\beta<m.
\end{equation}
Indeed, if the two smallest elements of $B$ have sum at least $m$, then the two largest elements $v<u$ of $B$ have $u+v>m$, since equality would give a forbidden sum equal to $m$. Reflecting sends them to the two smallest elements $m-u<m-v$ of $-B$, whose sum is $2m-u-v<m$.

We now split into two cases.

\smallskip
\noindent\textbf{Case 1: $\alpha+\beta\in C(B)$.}
Let
\[
  B'=B\setminus\{\alpha\}.
\]
Since $\alpha$ and $\beta$ are the two smallest elements of $B$, every low representation of $\alpha+\beta$ uses $\alpha$. Therefore, after deleting $\alpha$, the residue $\alpha+\beta$ no longer lies in $C(B')$. Deleting elements cannot create new collisions, so
\[
  |C(B')|\leq |C(B)|-1.
\]
Consequently
\[
  |B|-|C(B)|\leq |B'|-|C(B')|.
\]
The induction hypothesis applied to $B'$ gives the desired bound.

\smallskip
\noindent\textbf{Case 2: $\alpha+\beta\notin C(B)$.}
Work in the cyclic group $\Z/m\Z$. Consider the four sets
\[
  T_1=B,
  \qquad
  T_2=-B,
  \qquad
  T_3=(B-\alpha)\setminus\{0\},
  \qquad
  T_4=(\beta-B)\setminus\{0\},
\]
where these are viewed as subsets of $\Z/m\Z$. Their total size is
\begin{equation}\label{eq:four-total}
  |T_1|+|T_2|+|T_3|+|T_4|=4|B|-2.
\end{equation}
We record the required intersection bounds. First,
\[
  |T_1\cap T_2|\leq 1,
\]
and if $m$ is odd then $T_1\cap T_2=\varnothing$, since an element in this intersection gives two elements of $B$ whose sum is $m$, except possibly the single residue $m/2$. Also
\[
  |T_1\cap T_3|\leq 1,
\]
since $x\in T_1\cap T_3$ gives $x\in B$ and $x\equiv b-\alpha$ for some $b\in B$, so
\[
  \alpha+x\equiv b\pmod m;
\]
unless $x=\alpha$, this is a forbidden sum of two distinct elements of $B$ landing back in $B$. Similarly,
\[
  |T_1\cap T_4|\leq 2,
\]
since such an intersection gives $x+b\equiv \beta\pmod m$ with $x,b\in B$; unless $x=b$, this is forbidden, while $x=b$ forces $2x\equiv \beta\pmod m$, which has at most two solutions modulo $m$.

Next,
\[
  |T_2\cap T_3|\leq 1.
\]
Indeed, such an intersection gives $u+b\equiv \alpha\pmod m$ with $u,b\in B$. Unless $u=b$, this is forbidden, so $2u\equiv\alpha\pmod m$. This congruence has at most one solution in $B$: if $m$ is odd there is only one residue solution, while if $m$ is even then either there is no solution, or the two residue solutions are $\alpha/2$ and $\alpha/2+m/2$; the first is smaller than $\alpha$, so at most the second can lie in $B$.
Finally,
\[
  |T_2\cap T_4|\leq 1;
\]
indeed, such an intersection gives a congruence $u+\beta\equiv b\pmod m$ with $u,b\in B$, which is forbidden unless $u=\beta$.

It remains to bound $T_3\cap T_4$. Suppose
\[
  b-\alpha\equiv \beta-b'\pmod m
\]
with $b,b'\in B$. Then
\[
  b+b'\equiv \alpha+\beta\pmod m.
\]
Since $0<\alpha+\beta<m$, this means either
\[
  b+b'=\alpha+\beta
\]
or
\[
  b+b'=m+\alpha+\beta.
\]
The zero omissions in the definitions of $T_3$ and $T_4$ mean that $b\neq\alpha$ and $b'\neq\beta$. Thus in the first equality, since $\alpha,\beta$ are the two smallest elements of $B$, the only possible ordered contribution is $(b,b')=(\beta,\alpha)$. In the second equality, if $b\neq b'$, then $\alpha+\beta$ has both the low representation $\alpha+\beta$ and the high representation $b+b'-m$, contrary to $\alpha+\beta\notin C(B)$. If $b=b'$, then $2b=m+\alpha+\beta$, so there is at most one further possible residue. Hence
\[
  |T_3\cap T_4|\leq 2.
\]

Thus the total pairwise overlap is at most $8$ if $m$ is even and at most $7$ if $m$ is odd. Also $0$ lies in none of the sets $T_i$, so
\[
  \left|\bigcup_{i=1}^4T_i\right|\leq m-1.
\]
By the elementary union bound in the form
\[
  \sum_{i=1}^4 |T_i|
  \leq \left|\bigcup_{i=1}^4T_i\right|+
  \sum_{1\leq i<j\leq 4}|T_i\cap T_j|,
\]
we get $4|B|-2\leq m+7$ if $m$ is even and $4|B|-2\leq m+6$ if $m$ is odd. Thus $4|B|\leq m+9$ in the even case and $4|B|\leq m+8$ in the odd case. Since $4|B|$ is divisible by $4$, the even case also gives $4|B|\leq m+8$. Hence in all cases
\[
  |B|\leq \frac m4+2.
\]
Since $|C(B)|\geq 0$, this proves
\[
  |B|-|C(B)|\leq \frac m4+2.
\]
Together with the induction case, this proves the lemma.
\end{proof}

\section{Folding a triple-free set around a pivot}

Call a set $A\subseteq[1,N]$ \emph{triple-free} if it contains no pairwise-sum triple. Let $A$ be triple-free and let $h\in A$. Define
\[
  X=\{r:1\leq r<h,\ r\in A\},
\]
\[
  Y=\{r:1\leq r<h,\ h+r\leq N,\ h+r\in A\}.
\]
Put
\[
  B_h=X\cap Y,
  \qquad
  E=[1,h-1]\setminus(X\cup Y).
\]
Thus $B_h\subseteq\{1,\ldots,h-1\}$, and $r\in B_h$ means that both $r$ and $h+r$ lie in $A$.


\begin{lemma}[Folding lemma]\label{lem:folding}
One has
\[
  |X|+|Y|\leq \frac54h-|E\setminus C(B_h)|+1.
\]
In particular,
\[
  |X|+|Y|\leq \frac54h+1.
\]
\end{lemma}

\begin{proof}
For $h=1$ there is nothing to prove, so assume $h\geq 2$. Take distinct $x,y\in B_h$. Then
\[
  x,y,h,x+h,y+h\in A.
\]
If $x+y=h$, then $x,y,h$ form a forbidden triple. If $x+y\equiv r\pmod h$ for some $r\in B_h$, then either $x+y=r$ or $x+y=h+r$; in either case $x+y\in A$, and again $x,y,h$ form a forbidden triple. Hence $B_h$ satisfies the hypotheses of Lemma~\ref{lem:folded} modulo $h$.

We next show that
\begin{equation}\label{eq:collisions-empty}
  C(B_h)\subseteq E.
\end{equation}
If $r\in C(B_h)$, then $r$ has a low representation $r=x+y$ with distinct $x,y\in B_h$ and a high representation $h+r=u+v$ with distinct $u,v\in B_h$. If $r\in X$, then $r\in A$, and $x,y,h$ form a forbidden triple. Thus $r\notin X$. If $r\in Y$, then $h+r\in A$, and $u,v,h$ form a forbidden triple. Thus $r\notin Y$. Therefore $r\in E$.

Now
\[
  |X|+|Y|=|X\cup Y|+|X\cap Y|=(h-1-|E|)+|B_h|.
\]
Using \eqref{eq:collisions-empty} and Lemma~\ref{lem:folded},
\[
  |X|+|Y|=h-1+|B_h|-|C(B_h)|-|E\setminus C(B_h)|
  \leq \frac54h-|E\setminus C(B_h)|+1.
\]
The final assertion follows by discarding the nonnegative term $|E\setminus C(B_h)|$.
\end{proof}


\end{document}
